\documentclass{article}

\usepackage{amsmath}
\usepackage{amssymb}

\title{Übungsblatt 4}
\author{Pascal Diller, Timo Rieke}

\begin{document}

\maketitle

\section*{Aufgabe 1}
\subsection*{(i)}
Für $U_1$: \\
\[x + z = 0 \Longleftrightarrow x = -z\]
\[(x,y,z,w) = (-z, y, z, w) = y(0,1,0,0) + z(-1,0,1,0) + w(0,0,0,1)\]
Basis für $U_1$: $\{(0,1,0,0), (-1,0,1,0),(0,0,0,1)\}$ \\
Dimension von $U_1: \dim{U_1} = 3$ \\
\newline
Für $U_2$: \\
\[x+y-z = 0 \Longleftrightarrow x = z - y\]
\[y-w = 0 \Longleftrightarrow w = y\]
\[(x,y,z,w) = (z - y, y, z , y) = y(-1,1,0,1) + z(1,0,1,0)\]
Basis von $U_2$: $\{(-1, 1, 0, 1), (1,0,1,0)\}$ \\
Dimension von $U_2: \dim{U_2} = 2$

\subsection*{(ii)}
\[x+z = 0 \Longleftrightarrow x = -z\]
\[x+y-z = 0 = -z + y -z \Longleftrightarrow y = 2z\]
\[y - w = 0 \Longleftrightarrow y = w = 2z\]
\[(x,y,z,w) = (-z,2z,z,2z) = z(-1,2,1,2)\]
Basis: $\{(-1,2,1,2)\}$ \\
$\dim{U_1 \cap U_2} = 1$

\subsection*{(iii)}
\[\dim{(U_1 + U_2)} = \dim{(U_1)} + \dim{(U_2)} - \dim{(U_1 \cap U_2)} = 3 + 2 - 1 = 4\]

\section*{Aufgabe 2}
\subsection*{(ii)}
$M_n(\mathbb{R})$ ist die Menge aller $n \times n$-Matrizen mit Einträgen in $\mathbb{R}$, also:
\[\dim(M_n(\mathbb{R})) = n^2\]
Da $T$ bijektiv ist, müssen $\dim(M_n(\mathbb{R}))$ und $\dim(\mathbb{R}^m)$ gleich sein.
\[\dim(\mathbb{R}^m) = m = \dim(M_n(\mathbb{R})) = n^2\]
\[m = n^2\]

\section*{Aufgabe 3}
$v_1 \neq 0$, also linear unabhängig. \\
Prüfe $(v_1, v_2)$ auf lineare Unabhängigkeit:
\[a v_1 + b v_2 = 0 \Rightarrow a(1,2,1) + b(-1,3,-2) = (0,0,0)\]
\[\Rightarrow \begin{cases} a - b = 0 \\ 2a + 3b = 0 \\ a - 2b = 0 \end{cases} \Rightarrow a = b = 0\]
$\Rightarrow$ $(v_1, v_2)$ linear unabhängig. \\
Prüfe $(v_1, v_2, v_3)$:
\[a v_1 + b v_2 + c v_3 = 0\]
Beispiel: $2v_1 - 2v_2 - v_3 = 0$, also $v_3$ ist abhängig. \\
Prüfe $(v_1, v_2, v_4)$:
\[c_1 v_1 + c_2 v_2 = v_4 \Rightarrow c_1(1,2,1) + c_2(-1,3,-2) = (1,7,1)\]
\[\Rightarrow \begin{cases} c_1 - c_2 = 1 \\ 2c_1 + 3c_2 = 7 \\ c_1 - 2c_2 = 1 \end{cases} \Rightarrow \text{Widerspruch (3. Gleichung nicht erfüllt)}\]
$\Rightarrow$ $v_4$ ist nicht linear abhängig von $v_1, v_2$. \\
$\{v_1, v_2, v_4\}$ ist eine Basis von $\mathbb{R}^3$.

\section*{Aufgabe 4}
\[a w_1 + b w_2 = 0 \Rightarrow a(r v_1 + v_2) + b(v_1 + s v_2) = 0\]
\[\Rightarrow (a r + b) v_1 + (a + b s) v_2 = 0\]
Da $(v_1, v_2)$ eine Basis ist, folgt:
\[a r + b = 0\] 
\[ a + b s = 0\]
Das lineare Gleichungssystem hat genau dann nur die Lösung $a = b = 0$, wenn die Determinante ungleich null ist:
\[\det \begin{pmatrix} r & 1 \\ 1 & s \end {pmatrix}= rs - 1 \neq 0\]
$(w_1, w_2)$ ist genau dann eine Basis von $V$, wenn $rs \neq 1$.

\end{document}
